\subsection{Convex Functions and Jensen's Inequality}

Convexity is the geometric condition behind Jensen's inequality: averages move inside the argument, while supporting lines convert the nonlinear statement into a linear integral estimate.

\begin{definition} \label{def:convex-function-on-interval} \hlt{Convex Function on an Interval}\\
Let $I\subseteq\R$ be an interval and let $\varphi:I\to\R$.
\begin{enumerate}[label=(\roman*)]
\setlength{\itemsep}{0pt}
\item $\varphi$ is \hlt{convex} iff $\varphi(\lambda x_1+(1-\lambda)x_2)\leq\lambda\varphi(x_1)+(1-\lambda)\varphi(x_2)$ for all $x_1,x_2\in I$ and $\lambda\in[0,1]$.
\item Equivalently, for all $x_1<x<x_2$ in $I$, $\varphi(x)\leq\frac{x_2-x}{x_2-x_1}\varphi(x_1)+\frac{x-x_1}{x_2-x_1}\varphi(x_2)$.
\end{enumerate}
\hlt{Key Insight}: Graphs lie below enclosing chords.
\end{definition}

\begin{figure}[H]
  \centering
  \begin{tikzpicture}[
  scale=0.88,
  >=Stealth,
  every node/.style={font=\scriptsize},
  axis/.style={->, gray!70, line width=0.45pt},
  curve/.style={blue!70!black, line width=0.9pt},
  chord/.style={orange!80!black, line width=0.8pt},
  support/.style={green!45!black, line width=0.8pt},
  guide/.style={gray!45, dashed, line width=0.4pt},
  point/.style={circle, fill=gray!85, inner sep=1.25pt}
]
  \draw[axis] (-0.12,0) -- (6.35,0) node[right] {$x$};
  \draw[axis] (0,-0.08) -- (0,2.55) node[above] {$y$};

  \coordinate (xone) at (1.10,1.422);
  \coordinate (alpha) at (3.20,0.54);
  \coordinate (xtwo) at (5.30,1.422);

  \draw[support] (0.62,0.54) -- (5.86,0.54);
  \draw[chord] (xone) -- (xtwo);
  \draw[curve, domain=0.58:5.82, samples=120]
    plot (\x,{0.20*(\x-3.2)*(\x-3.2)+0.54});

  \draw[guide] (1.10,0) -- (xone);
  \draw[guide] (3.20,0) -- (alpha);
  \draw[guide] (5.30,0) -- (xtwo);

  \node[below] at (1.10,0) {$x_1$};
  \node[below] at (3.20,0) {$\alpha$};
  \node[below] at (5.30,0) {$x_2$};

  \node[point] at (xone) {};
  \node[point] at (alpha) {};
  \node[point] at (xtwo) {};
\end{tikzpicture}

  \caption{Color code: \textcolor{blue!70!black}{$y=\varphi(x)$} lies below the \textcolor{orange!80!black}{chord} through $x_1,x_2$ and above the \textcolor{green!45!black}{supporting line} at $\alpha$.}
  \label{fig:convex-chord-supporting-line}
\end{figure}

\begin{lemma} \label{lem:convex-chordal-slope-lemma} \hlt{Chordal Slope Lemma}\\
Let $I\subseteq\R$ be an interval, let $\varphi:I\to\R$ be convex, and let $x_1,x,x_2\in I$ satisfy $x_1<x<x_2$.\\
Then $\frac{\varphi(x)-\varphi(x_1)}{x-x_1}\leq\frac{\varphi(x_2)-\varphi(x_1)}{x_2-x_1}\leq\frac{\varphi(x_2)-\varphi(x)}{x_2-x}$.
\begin{proof} Since $x=\frac{x_2-x}{x_2-x_1}x_1+\frac{x-x_1}{x_2-x_1}x_2$, convexity gives $(x_2-x_1)\varphi(x)\leq(x_2-x)\varphi(x_1)+(x-x_1)\varphi(x_2)$.\\
Rearrange, $(x_2-x_1)(\varphi(x)-\varphi(x_1))\leq(x-x_1)(\varphi(x_2)-\varphi(x_1))$; $(x_2-x)(\varphi(x_2)-\varphi(x_1))\leq(x_2-x_1)(\varphi(x_2)-\varphi(x))$.\\
Divide by $(x-x_1)(x_2-x_1)>0$ and $(x_2-x)(x_2-x_1)>0$, respectively.\\
\hlt{Structural Commentary}: One chord inequality, two rearrangements.\\
\hlt{Key Insight}: Convex secant slopes are nondecreasing.
\end{proof}
\end{lemma}

\begin{proposition} \label{prop:increasing-derivative-convex-function} \hlt{Increasing Derivative Gives Convexity}\\
Let $I\subseteq\R$ be open and let $\varphi:I\to\R$ be differentiable. If $\varphi'$ is increasing on $I$, then $\varphi$ is convex on $I$.\\
In particular, this holds if $\varphi$ is twice differentiable and $\varphi''\geq0$ on $I$.
\begin{proof} For $a<b$ in $I$, set $S(a,b)\defeq(\varphi(b)-\varphi(a))/(b-a)$. Let $u<v$ and $0<\lambda<1$, and put $x=\lambda u+(1-\lambda)v$.\\
By MVT, $S(u,x)=\varphi'(c_1)$, $S(x,v)=\varphi'(c_2)$ for some $c_1\in(u,x)$, $c_2\in(x,v)$. As $c_1<c_2$, $S(u,x)\leq S(x,v)$.\\
Clearing denominators, $\varphi(x)\leq\frac{v-x}{v-u}\varphi(u)+\frac{x-u}{v-u}\varphi(v)=\lambda\varphi(u)+(1-\lambda)\varphi(v)$; the cases $\lambda=0,1$ are equality.\\
If $\varphi''\geq0$, for $a<b$, MVT gives $\varphi'(b)-\varphi'(a)=\varphi''(d)(b-a)\geq0$ for some $d\in(a,b)$, so $\varphi'$ is increasing.\\
\hlt{Structural Commentary}: MVT orders secant slopes.\\
\hlt{Key Insight}: Monotone slopes give convexity.
\end{proof}
\end{proposition}

\begin{example} \label{ex:basic-convex-functions-second-derivative} \hlt{Basic Convex Functions}\\
By Proposition \ref{prop:increasing-derivative-convex-function}, $x\mapsto x^p$ for $p\geq1$, $x\mapsto e^{ax}$, and $x\mapsto\log(1/x)$ are convex on $(0,+\infty)$, $\R$, and $(0,+\infty)$.
\end{example}

\begin{proposition} \label{prop:finite-jensen-inequality} \hlt{Finite Jensen Inequality}\\
Let $I\subseteq\R$ be an interval, $\varphi:I\to\R$ be convex, $x_1,\ldots,x_n\in I$, $\lambda_1,\ldots,\lambda_n\geq0$ with $n\in\N$, $\sum_{k=1}^{n}\lambda_k=1$.\\
Set $\bar{x}\defeq\sum_{k=1}^{n}\lambda_kx_k$. Then $\bar{x}\in I$ and $\varphi(\bar{x})\leq\sum_{k=1}^{n}\lambda_k\varphi(x_k)$.
\begin{proof} Induct on $m=\#\{k:\lambda_k>0\}$. If $m=1$, equality holds.\\
For $m\geq2$, relabel so $\lambda_1,\ldots,\lambda_m>0$ and $\lambda_k=0$ for $k>m$. Let $\Lambda=1-\lambda_m$ and $y=\Lambda^{-1}\sum_{k=1}^{m-1}\lambda_kx_k$.\\
By induction, $y\in I$ and $\varphi(y)\leq\Lambda^{-1}\sum_{k=1}^{m-1}\lambda_k\varphi(x_k)$.\\
Since $\bar{x}=\Lambda y+\lambda_mx_m$, convexity gives $\bar{x}\in I$ and $\varphi(\bar{x})\leq\Lambda\varphi(y)+\lambda_m\varphi(x_m)\leq\sum_{k=1}^{n}\lambda_k\varphi(x_k)$.\\
\hlt{Structural Commentary}: Induct by peeling off one weight.\\
\hlt{Key Insight}: Jensen iterates convexity.
\end{proof}
\end{proposition}

\begin{lemma} \label{lem:convex-one-sided-derivatives-supporting-lines} \hlt{One-Sided Slopes and Supporting Lines}\\
Let $I\subseteq\R$ be open and let $\varphi:I\to\R$ be convex.
\begin{enumerate}[label=(\roman*)]
\setlength{\itemsep}{0pt}
\item For every $x\in I$, $\varphi'(x-)$ and $\varphi'(x+)$ exist and are finite.
\item If $u<v$ in $I$, then $\varphi'(u-)\leq\varphi'(u+)\leq\frac{\varphi(v)-\varphi(u)}{v-u}\leq\varphi'(v-)\leq\varphi'(v+)$.
\item If $\alpha\in I$ and $\varphi'(\alpha-)\leq m_0\leq\varphi'(\alpha+)$, then $\varphi(t)\geq\varphi(\alpha)+m_0(t-\alpha)$ for all $t\in I$.
\end{enumerate}
\begin{proof}
For $a<b$ in $I$, write $S(a,b)\defeq(\varphi(b)-\varphi(a))/(b-a)$.
\begin{enumerate}[label=(\roman*)]
\setlength{\itemsep}{0pt}
\item Fix $x\in I$ and choose $a<x<b$ in $I$.\\
If $a<r<s<x$, Lemma \ref{lem:convex-chordal-slope-lemma} gives $S(a,x)\leq S(r,x)\leq S(s,x)\leq S(x,b)$, so $\lim_{r\uparrow x}S(r,x)$ exists in $\R$.\\
If $x<r<s<b$, the same lemma gives $S(a,x)\leq S(x,r)\leq S(x,s)\leq S(x,b)$, so $\lim_{r\downarrow x}S(x,r)$ exists in $\R$.
\item Let $u<v$ in $I$. For $r<u<s<v$ and $u<q<v<t$, the chordal slope inequalities give\\
$S(r,u)\leq S(u,s)\leq S(u,v)\leq S(q,v)\leq S(v,t)$. Letting $r\uparrow u$, $s\downarrow u$, $q\uparrow v$, and $t\downarrow v$ gives asserted chain.
\item If $t=\alpha$, equality holds. If $t>\alpha$, then (ii) gives $m_0\leq\varphi'(\alpha+)\leq S(\alpha,t)$, hence $\varphi(t)\geq\varphi(\alpha)+m_0(t-\alpha)$.\\
If $t<\alpha$, then (ii) gives $S(t,\alpha)\leq\varphi'(\alpha-)\leq m_0$, which rearranges to the same estimate.
\end{enumerate}
\hlt{Structural Commentary}: Monotone secant slopes have finite one-sided limits.\\
\hlt{Key Insight}: One-sided slope bounds produce supporting lines.
\end{proof}
\end{lemma}

\begin{theorem} \label{thm:jensen-inequality-lebesgue-unit-interval} \hlt{Jensen's Inequality on $[0,1]$}\\
Let $m$ be Lebesgue measure on $[0,1]$, let $\varphi:\R\to\R$ be convex, and let $f,\varphi\circ f\in L^1(m)$.\\
Then $\varphi\left(\int f\,dm\right)\leq\int\varphi(f)\,dm$.
\begin{proof}
Set $\alpha\defeq\int f\,dm$. By Lemma \ref{lem:convex-one-sided-derivatives-supporting-lines}, choose $m_0\in[\varphi'(\alpha-),\varphi'(\alpha+)]$; $\varphi(t)\geq\varphi(\alpha)+m_0(t-\alpha)$ $\forall t\in\R$.\\
Substituting $t=f$ and integrating gives $\int\varphi(f)\,dm\geq\varphi(\alpha)+m_0(\int f\,dm-\alpha)=\varphi(\alpha)$.\\
\hlt{Structural Commentary}: Support at the mean.\\
\hlt{Key Insight}: The linear term vanishes.
\end{proof}
\end{theorem}

\begin{remark} \label{rem:jensen-integrability-growth-condition} \hlt{Integrability in Jensen}\\
Linear growth, $\abs{\varphi(t)}\leq c_1+c_2\abs t$, makes $\varphi\circ f\in L^1([0,1])$ by Proposition \ref{prop:lp-composition-growth-and-continuity}; if $\varphi\circ f\geq0$, Jensen's right side may be extended-valued.
\end{remark}
